\section{From homogeneous jets to nefness}\label{sec:nef}

The preceding sections produce an injective jet test at some tuple for
each pair of integers $k,m$ above the critical slope. We first make all
these tests hold at one very general tuple. A separate surface argument
then converts them into inequalities for every curve, including curves
whose multiplicities at the points are unequal.

\begin{lemma}[The open locus of an injective jet test]\label{lem:jet-open}
Let $S$ be a smooth integral complex projective surface, let $L$ be a line
bundle, and let $r,k,m$ be positive integers. In the configuration space
\[
 U_r=\{(p_1,\ldots,p_r)\in S^r:p_i\ne p_j\text{ for }i\ne j\},
\]
the tuples at which the order-$m$ jet test on $H^0(S,L^{\otimes k})$
is injective form a Zariski-open subset. If each test in a countable
collection with fixed $S,L,r$ is injective at some tuple, then all those tests are
simultaneously injective at very general tuples. Their common locus is
nonempty.
\end{lemma}

\begin{proof}
Let $q_1,q_2:S\times S\to S$ be the projections and let
$\mathcal I_\Delta$ be the ideal of the diagonal. The bundle of jets
of orders less than $m$ is
\[
 \mathcal J_{k,m}
 =q_{1*}\!\left(q_2^*L^{\otimes k}\otimes
       \mathcal O_{S\times S}/\mathcal I_\Delta^m\right).
\]
This is the algebraic principal-parts construction
\cite[Tags 0G3P and 0638]{Stacks}; it is locally free of rank
$\binom{m+1}{2}$.
Indeed, the diagonal is a regular embedding with conormal bundle
$\Omega_S^1$; the successive quotients of the displayed sheaf, pushed
forward to $S$, are
$L^{\otimes k}\otimes\operatorname{Sym}^j\Omega_S^1$ for
$0\le j<m$. They are locally free of ranks $j+1$.
The thickened diagonal is finite over $S$, so these pushforwards are
exact and give the asserted filtration. The fiber at $p$ is
$(L^{\otimes k})_p/\mathfrak m_p^m(L^{\otimes k})_p$.

Writing $\operatorname{pr}_i:U_r\to S$ for the coordinate projections,
consider the algebraic bundle map
\[
 H^0(S,L^{\otimes k})\otimes\mathcal O_{U_r}
 \longrightarrow\bigoplus_{i=1}^r\operatorname{pr}_i^*\mathcal J_{k,m}.
\]
Its fiber kernel consists exactly of sections vanishing to order at
least $m$ at every point. If $n=h^0(S,L^{\otimes k})>0$, failure of
injectivity is cut out locally by the $n\times n$ minors. If $n=0$,
the map is injective everywhere. In either case the injective locus
is Zariski open, and a single injective tuple makes it nonempty.
The analytic and algebraic finite jets agree at complex points:
in local coordinates both are the same Taylor polynomial modulo terms
of total degree at least $m$. Thus a tuple obtained by the analytic
transfers of the preceding sections is a valid witness here.

The variety $U_r$ is smooth and irreducible. A proper algebraic closed
subset has empty interior in its complex manifold of points. For a
countable collection of nonempty injective loci, their complements are
therefore closed nowhere-dense subsets of $U_r(\C)$. This locally
compact Hausdorff space is a Baire space, so their common complement
is nonempty, indeed dense. By construction it is the complement of a
countable union of proper Zariski-closed subsets of $U_r$, as required
for very general tuples.
\end{proof}

We next record the precise Riemann--Roch argument that will turn a
negative curve into an effective divisor with equal prescribed
multiplicities. The divisor being made effective need not be nef.

\begin{lemma}[Effective multiples of a positive-square divisor]
\label{lem:positive-square}
Let $Y$ be a smooth integral complex projective surface, let $A$ be a
nef real divisor class, and let $D$ be a rational Cartier divisor.
If $D^2>0$ and $A\cdot D>0$, then
$H^0(Y,\mathcal O_Y(kD))\ne0$ for every sufficiently large positive
integer $k$ divisible by the denominators of $D$.
\end{lemma}

\begin{proof}
Choose a positive integer $q$ such that $qD$ is Cartier and let
$k$ run through positive multiples of $q$. By Serre duality \cite[Tag 0FVV, Lemma 48.27.1(6)--(7)]{Stacks},
\[
 h^2(Y,\mathcal O_Y(kD))
   =h^0(Y,\mathcal O_Y(K_Y-kD)).
\]
For sufficiently large $k$,
\[
 A\cdot(K_Y-kD)=A\cdot K_Y-kA\cdot D<0.
\]
A divisor with a nonzero section is linearly equivalent to an
effective divisor, which has nonnegative intersection with the nef
class $A$. Hence the displayed $h^2$ vanishes for all such $k$.
Surface Riemann--Roch \cite[Tag 02UO]{Stacks} now gives
\[
 \begin{split}
 h^0(Y,\mathcal O_Y(kD))
 &\ge \chi(Y,\mathcal O_Y(kD))\\
 &=\frac{k^2D^2-kD\cdot K_Y}{2}+\chi(Y,\mathcal O_Y)>0
 \end{split}
\]
for sufficiently large $k$, since $D^2>0$.
\end{proof}

\begin{proposition}[Homogeneous tests imply nefness]
\label{prop:homogeneous-nef}
Let $S$ be a smooth integral complex projective surface, let $L$ be
ample, and let $p_1,\ldots,p_r$ be distinct points. Put
$H=L^2$ and $w=\sqrt{H/r}$. Suppose that, for all positive integers
$k,m$ with $m>kw$, no nonzero section of $L^{\otimes k}$ vanishes to
order at least $m$ at every $p_i$. On the blow-up
$\pi:Y\to S$ of these points, with exceptional curves $E_i$, the real
divisor class
\[
 D_0=\pi^*L-w\sum_{i=1}^rE_i
\]
is nef.
\end{proposition}

\begin{proof}
Choose a Cartier divisor representing $L$ and retain the notation $L$
for it. Write $A=\pi^*L$ and $E=\sum_iE_i$. The class $A$ is nef and
\[
 A^2=H,\qquad A\cdot E_i=0,\qquad
 E_i\cdot E_j=-\delta_{ij},\qquad D_0^2=0.
\]
If $D_0$ were not nef, there would be an integral curve $C\subset Y$
with $D_0\cdot C<0$. Since $D_0\cdot E_i=w>0$, this curve is not
exceptional, and $A\cdot C=L\cdot\pi_*C>0$.
For sufficiently small positive rational $\delta$, both
\[
 (D_0-\delta C)^2=-2\delta D_0\cdot C+\delta^2C^2>0,
 \qquad
 A\cdot(D_0-\delta C)=H-\delta A\cdot C>0
\]
hold. Fix such a $\delta$. By continuity, there is a rational number
$s>w$, sufficiently close to $w$, such that
\[
 D=A-sE-\delta C
 \quad\text{satisfies}\quad D^2>0,\qquad A\cdot D>0.
\]
Here $A\cdot D=H-\delta A\cdot C$ is independent of $s$.

Choose a sufficiently large positive integer $k$ divisible by the
denominators of both $s$ and $\delta$.
Lemma~\ref{lem:positive-square} supplies an effective divisor
$F\sim kD$. Since $k\delta$ and $m=ks$ are positive integers,
\[
 F+k\delta C\sim kA-mE
\]
is effective. The blow-up identity \cite[Tag 0AGP, Lemma 54.3.4(4)]{Stacks}
\[
 \pi_*\mathcal O_Y(-mE)
   =\bigcap_{i=1}^r\mathcal I_{p_i}^{\,m}
\]
and the projection formula give a nonzero section of
$L^{\otimes k}$ vanishing to order at least $m$ at every $p_i$.
To recall the local content of this identity, sections of a pullback
line bundle descend across the missing points because $S$ is normal,
and the exceptional order of a local equation on the blow-up of a
smooth point is its ordinary multiplicity. Thus the condition at each
$E_i$ is precisely membership in $\mathcal I_{p_i}^{\,m}$.
But $m=ks>kw$, contradicting the hypothesis. Therefore $D_0$ is nef.
\end{proof}

For completeness, Riemann--Roch and a count of jet conditions also give
the universal upper bound for the Seshadri constant.

\begin{lemma}[The square upper bound]\label{lem:sesha-upper}
For every ample line bundle $L$ on a smooth integral complex projective
surface $S$ and every tuple of $r$ distinct points,
\[
 \varepsilon(S,L;p_1,\ldots,p_r)\le\sqrt{L^2/r}.
\]
\end{lemma}

\begin{proof}
Put $H=L^2$, let $\pi:Y\to S$ be the blow-up, and write
$A=\pi^*L$ and $E=\sum_iE_i$. Suppose $A-\lambda E$ is nef for
some $\lambda\ge0$. Fix $0<\alpha<\sqrt{H/r}$ and set
$m_k=\lceil k\alpha\rceil$. The proof of
Lemma~\ref{lem:positive-square}, applied on $S$ with $A=D=L$,
gives
\[
 h^0(S,L^{\otimes k})\ge \frac{H}{2}k^2+O(k).
\]
The dimension of the target of the order-$m_k$ jet evaluation at
the $r$ points is
\[
 r\binom{m_k+1}{2}=\frac{r\alpha^2}{2}k^2+O(k).
\]
Since $r\alpha^2<H$, evaluation has a nonzero kernel for all
sufficiently large $k$. A section in that kernel gives an effective
divisor $F_k\sim kA-m_kE$ on $Y$. Nefness implies
\[
 0\le(A-\lambda E)\cdot F_k=kH-\lambda r m_k.
\]
Letting $k\to\infty$ gives $\lambda\le H/(r\alpha)$, and then
letting $\alpha\uparrow\sqrt{H/r}$ gives
$\lambda\le\sqrt{H/r}$. Taking the supremum over the nef
coefficients $\lambda$ proves the claim.
\end{proof}

\begin{proof}[Proof of Theorem~\ref{thm:main}]
Choose the nodal divisor and parameters in
Section~\ref{sec:node}. Corollary~\ref{cor:torus-test} supplies an
integer $r_0$, depending only on these choices and hence on $(S,L)$,
with the following property. For every integer $r\ge r_0$, there is
a parameter $t>T$ and an $r$-point torus tuple such that the jet test
on $\mathcal M(kP_t)$ is injective for every pair of positive integers $k,m$
with $m>k\sqrt{H/r}$.

Fix any such integer $r$. For each of these pairs,
Corollary~\ref{cor:surface-transfer} gives an injective jet test on
$H^0(S,L^{\otimes k})$ at some tuple of $r$ distinct points of $S$.
The tuples may depend on $k,m$. Lemma~\ref{lem:jet-open} makes all
these countably many tests hold simultaneously at very general
tuples. Proposition~\ref{prop:homogeneous-nef} proves that the class
$\pi^*L-\sqrt{H/r}\sum_iE_i$ is nef at every such tuple.
Combining this lower bound for the Seshadri constant with
Lemma~\ref{lem:sesha-upper} proves the required equality.
The integer $r$ was arbitrary subject to $r\ge r_0$, and $r_0$
was fixed before $k,m$ were chosen. This proves the assertion for
every sufficiently large integer $r$.
\end{proof}

\subsection{Extension of the ground field}

\begin{corollary}\label{cor:fields}
Let $K$ be an uncountable algebraically closed field of characteristic
zero, let $S$ be a smooth integral projective surface over $K$, and let
$L$ be an ample line bundle on $S$. There is an integer $r_0=r_0(S,L)$
such that, for every $r\ge r_0$, outside a countable union of proper
Zariski-closed subsets of $U_r$, with nonempty complement, one has
\[
 \varepsilon(S,L;p_1,\ldots,p_r)=\sqrt{L^2/r}.
\]
\end{corollary}
\begin{proof}
Descend the projective equations and the finite gluing data for $L$ to
a subfield $F\subset K$ finitely generated over $\mathbb Q$. This gives
a smooth geometrically integral projective surface $S_0/F$ and an ample
line bundle $L_0$ whose base change is $(S,L)$; these properties descend
through field extension, and the finite-data descent and ampleness
criteria are recorded in \cite[Tags 01ZM and 0D3C]{Stacks}.
The field $F$ is countable and embeds into $\C$. Fix such an embedding
and choose $r_0$ from Theorem~\ref{thm:main} for $(S_{0,\C},L_{0,\C})$.
The intersection number $H=L_0^2=L^2$ is unchanged by field extension.

Fix $r\ge r_0$ and put $w=\sqrt{H/r}$. For every pair of positive
integers $a,m$ with $m>aw$, form on $U_{r,0}=U_r(S_0)$ the evaluation map
\[
 H^0(S_0,L_0^{\otimes a})\otimes\mathcal O_{U_{r,0}}
 \longrightarrow
 \bigoplus_{i=1}^r\operatorname{pr}_i^*\mathcal J_{a,m}
\]
using the principal-parts construction in the proof of
Lemma~\ref{lem:jet-open}. Its noninjectivity locus $B_{a,m}$ is closed,
cut out by maximal minors, and is empty if the source space is zero.
Principal parts and these minors commute with field extension, as does
$H^0$ by flat base change \cite[Tag 02KH]{Stacks}.
Each $B_{a,m}$ is proper over $\C$: at a tuple where
$A-wE$ is nef, with $A=\pi^*L_{0,\C}$ and $E=\sum_iE_i$, a section in
the evaluation kernel would give an effective divisor of class $aA-mE$,
whose intersection with $A-wE$ is $aH-rmw<0$.
Thus $B_{a,m}$ is proper over $F$ and remains proper over $K$.
The single threshold $r_0$ works for all these pairs.

Let $Z_r=\bigcup_{a,m:\,m>aw}B_{a,m,K}$. This is a countable union of
proper closed subsets. Its complement contains a $K$-point: the
function field $F(U_{r,0})$ has transcendence degree $2r$ over $F$ and
embeds over $F$ into $K$. Indeed, choose $2r$ algebraically independent
elements of $K$ over the countable field $F$ and extend the resulting
embedding of a rational function field across the finite algebraic
extension, using that $K$ is algebraically closed. The resulting
$K$-point maps to the generic point of $U_{r,0}$ and therefore avoids
all the $B_{a,m}$.

At every tuple outside $Z_r$, all the homogeneous jet tests are
injective. The proofs of Lemma~\ref{lem:positive-square},
Proposition~\ref{prop:homogeneous-nef}, and Lemma~\ref{lem:sesha-upper}
use only surface Riemann--Roch, Serre duality, and the blow-up ideal
identity, and are valid over $K$. They give nefness of
$\pi^*L-w\sum_iE_i$ and the matching upper bound, respectively.
\end{proof}
